\section{Conclusion}\label{sec:conclusion}

Is innovation too fast? The model's answer is that the question is usually asked about
the wrong margin. When papers arrive faster than they can be read, researchers do
not keep piling into the same topics until duplication explodes. They spread
out, and the overlap rate stays where entry pins it. The arXiv record bears this
out: the volume of papers more than doubled while close overlap barely moved.
The costs of speed lie elsewhere. The gains from building on others' work are competed away.
Incumbents keep rents behind a wall of reading that newcomers must climb.
Returns rise with density even though topics are too crowded at the margin.
And topics whose literature grows faster than newcomers can catch up overshoot.
Literally, they are too fast.

Several simplifications deserve scrutiny in future work. Credit is measured by
citations, which also reward audience size. This biases measured returns toward
dense topics, so the overcrowding we find is, if anything, understated.
Newcomers value topics at their current state, and diffusion rates are
exogenous. Forward-looking entry into a growing topic would anticipate
crowding and dampen overshooting, while information cascades would amplify it.
Topics are identical up to fertility. Their spatial arrangement in an area's
idea space appears only in prediction P7. And the frontier is a reduced-form
lottery. Endogenising it, for instance by letting pioneers choose how far from
existing topics to locate, is the natural next step. It is also the margin
that the policy conclusion below hinges on.

The policy implications follow from where the rents go. Anything that improves the payoff of
working inside existing topics is competed away. That includes better reading
tools, search engines, summarisers and caps on writing. Anything that improves
the frontier is not. Writing tools scale science because they scale the frontier
too. Reading tools crowd it unless they also make new topics easier to open. To
date, artificial intelligence has changed science mainly by lowering the cost of
entry, and the new entrants have spread out. The model predicts what will happen when tools that read for
us become widespread: denser topics, fewer new ones, more hot topics that
overshoot, and little gain in output per researcher. Unless, that is, the same
tools are pointed at the frontier.
